% ---------------------------------------------------------------- preambulo
% Comum aos dois artigos (2D e 3D).  Mantido num arquivo so' para que uma
% mudanca de estilo nao precise ser feita duas vezes e divergir na segunda.
\usepackage[T1]{fontenc}
\usepackage[utf8]{inputenc}
\usepackage[english]{babel}
\usepackage{lmodern}
\usepackage{microtype}
\usepackage{amsmath,amssymb}
\usepackage{booktabs}
\usepackage{adjustbox}
\usepackage{siunitx}
\usepackage{graphicx}
\usepackage{pgfplots}
\usepackage{pgfplotstable}
\usepackage{tikz}
\usepackage{xcolor}
\usepackage{capt-of}
\pgfplotsset{compat=1.18}

% OS NUMEROS VIVEM EM dados/*.dat e sao lidos por pgfplotstable, nunca
% digitados no texto.  Regerar uma medida troca o .dat e o texto acompanha,
% sem chance de a tabela e a prosa discordarem.
\newcommand{\dat}[1]{dados/#1}

% \status{...}: o que a secao tem de respaldo medido.  Definir \nostatus antes
% do \begin{document} faz os blocos sumirem na versao de submissao.
\newcommand{\status}[1]{%
  \ifdefined\nostatus\else
    \begingroup
    \par\medskip\noindent
    \colorbox{yellow!18}{\begin{minipage}{0.97\linewidth}
      \small\textbf{Status of this section.} #1
    \end{minipage}}\par\medskip
    \endgroup
  \fi}

% \missing{...}: uma lacuna pontual, visivel de proposito.
\newcommand{\missing}[1]{\textcolor{red!70!black}{\textbf{[missing:} #1\textbf{]}}}

% \caveat NAO e' o mesmo que \missing.  \missing marca o que falta MEDIR;
% \caveat marca o que foi medido mas com uma ressalva que muda a leitura.
% Confundir os dois esconderia um resultado-com-reserva entre pendencias.
\newcommand{\caveat}[1]{\par\smallskip\noindent\textcolor{orange!60!black}{%
  \textbf{[caveat:} #1\textbf{]}}\par\smallskip}
